\section{Discussion}
\label{sec:discussion}

\subsection{Why Explanation Signals Can Complement Uncertainty}

The active-learning result suggests that output ambiguity and explanation
geometry capture related but non-identical information. Predictive uncertainty
depends on one scalar. Directional entropy distinguishes one-sided support from
strong opposing contributions, while SHAP variance distinguishes diffuse from
concentrated attributions. Their usefulness is plausible in a PE setting where
benign-looking structural metadata may coexist with suspicious imports, sections,
or byte patterns. Two samples near $p=0.5$ need not be equally informative if one
contains little evidence and the other contains a conflict among large effects.

The ablation results also caution against assigning each signal an unconditional
benefit. $U+H$ had a favorable mean FPR but lower TPR and AUC than $\xfull$;
$U+V$ had relatively strong AP but unstable FPR. The full weighted score performed
best on mean AUC and AP, yet the weights were evaluated on one local dataset and
were not optimized through an independent outer validation study. A larger study
should compare fixed preregistered weights, validation-only tuning, rank
aggregation, and diversity-aware batch selection. It should also test whether
acquisition becomes redundant when many high-scoring pool samples are near
duplicates.

The current result demonstrates ranking improvement at equal training-label
counts. It does not demonstrate label savings, because a target validation
performance was not declared before the run, and it does not demonstrate labor
savings, because expert annotation was simulated by label revelation. A
deployment study would measure time per label, adjudication, abstention, and
mislabeling, ideally with multiple analysts.

\subsection{What the Fidelity Experiment Establishes}

The deletion control answers a model-behavior question: if the coordinates that
TreeSHAP ranks highest are replaced by a training statistic, does predictive
ranking degrade more than under random replacement? Across all three budgets the
answer is yes, and the effect grows with $k$. The random baseline is essential.
Without it, an AUC decrease could be attributed merely to modifying input
coordinates. The small random drops show that the same nominal number of
uniformly selected coordinates usually has little effect on this model.

The experiment does not establish executable-level fidelity. A PE representation
contains dependent quantities: sizes, offsets, counts, hashes, and directory
metadata may obey structural constraints. Coordinate-wise mean replacement can
violate those constraints or create a vector far from the data manifold. Even
equal $k$ does not ensure equal perturbation magnitude; some selected values may
already equal their means, while a top-SHAP coordinate may be far from baseline.
The artifact records the actual number of changed coordinates and probability
change, but a stronger protocol would mask semantically coupled groups, quantify
out-of-distribution distance, and where possible reconstruct and validate a PE
file.

The approximately 630-fold ratio at $k=25$ should not be marketed as a universal
improvement factor. It divides by a random-control mean close to zero and can vary
substantially with baseline, $k$, model, split, and random seed. The more stable
statement is that the absolute top-SHAP AUC drop (0.1616) is materially larger
than the random mean (0.000257) under the frozen protocol, with a conditional
paired-bootstrap interval [0.1222, 0.2078].

EC and CF add geometric context but do not certify an explanation. EC depends on
the chosen feature standardization, cosine similarity, and Pearson comparison.
CF uses test labels and is sensitive to centroid structure and outliers. Because
FF, EC, and CF measure different constructs and do not naturally share a bounded
scale, an aggregate OFS would require a preregistered normalization, justified
weights, sensitivity analysis, and ideally correlation with external expert or
task outcomes. Reporting the components separately is scientifically stronger
than selecting a convenient weighted average after observing them.

\subsection{Role of PID in the Multi-Stage System}

The PID proposal is motivated by sequential disagreement, but its value cannot be
inferred from the controller form. A fixed average may be difficult to beat if
the two sources are stable and calibrated. A proportional rule may be sufficient
if sessions are short. Integral state can accumulate bias, and a derivative term
can amplify noisy score changes. These possibilities make P/PI/PID ablations
necessary rather than ceremonial.

A proper dynamic study must distinguish the value of the source from the value
of fusion. Static-only versus Dynamic-only quantifies individual ranking; Average
and validation-selected Fixed establish simple fusion baselines; P, PI, and PID
test whether state and temporal change contribute. The same paired inputs must be
used for every rule. If PID does not improve the baselines, the correct outcome
is to report that result and prefer the simpler model.

The claim about ``sleep evasion'' also requires care. PID cannot create behavior
that a sandbox never observes. A longer or adaptive observation policy may expose
delayed behavior; the controller can only reweight scores after dynamic evidence
arrives. Likewise, sub-millisecond arithmetic replay is not evidence of a
15-ms end-to-end decision if process startup and observation take seconds. A
future latency table should separately report static extraction, static inference,
queueing, environment startup, dynamic observation, dynamic scoring, fusion, and
end-to-end time, including median and tail quantiles.

\subsection{Operational Implications}

The framework is best understood as an allocation system. Static analysis
provides broad, early coverage. A gate allocates expensive observation to selected
cases. Active learning allocates expert labels to selected training samples.
Fidelity testing allocates trust by requiring explanations to pass controlled
behavioral checks. These allocation decisions operate at different time scales
and should have independent service-level metrics.

For a security operations center, the most immediate verified use is model
maintenance: rank uncertain pool items using $\xfull$, present the selected PE
records and evidence to analysts, and retain the queries and subsequent model
versions. The fidelity protocol can serve as a release-time diagnostic for each
candidate model. The dynamic/PID path should remain disabled as a performance
claim until collection and validation gates pass; a software implementation can
exist before its operational efficacy has been established.

Reproducibility artifacts support audit and incident review. A model checkpoint,
query list, split manifest, prediction table, and hash make it possible to answer
why a reported metric changed. They do not replace governance requirements such
as sample-handling authorization, privacy controls, malware containment, or
documented label definitions. High-quality operational evidence combines both
computational traceability and trustworthy data lineage.
